%    Primer capítulo: Introducción a los lenguajes de programación.
%    Ejercicios por Barón M. Miguel.
%    Teoría por Javier Enríquez Mendoza.
%    Empezado el 17/10/22
%    Concluido el 20/10/22

%    Sección 1: Composición de Lenguajes.
\begin{figure}[h]
    \centerline{\includegraphics[scale=.33]{assets/01_gatito_introduccion.PNG}}       
\end{figure}



\bigskip

En el contexto de la computación, los lenguajes de programación fungen como una interfaz entre programador y procesador para establecer una comunicación sobre lo que queremos que la computadora ejecute por nosotros.\\\\
Diferentes paradigmas de lenguajes de programación, estilos, reglas y convenciones se han adoptado en las últimas décadas para ajustarse al mejor planteamiento posible de un problema, mantener un estándar declarativo, un nombramiento homogéneo de variables y métodos, o un estilo idéntico compartido entre los diferentes programas que se agrupan en estas categorías.\\\\
Iniciaremos el estudio de este manual planteando los conceptos claves que nos permitirán definir, clasificar y analizar a los lenguajes de programación que existen en el contexto de las ciencias de la computación. Estos conceptos nos permitirán definir nuestro propio lenguaje de programación desde cero, creando las estructuras y reglas necesarias para cada tipo, variable e instrucción del mismo.\\\\
Por último analizaremos sus propiedades y estudiaremos sus características más importantes para poder dar una implementación robusta aplicando las definiciones que iremos presentando en el desarrollo de esta introducción.\\

\subsection*{Objetivo}
    Comenzar el estudio formal de los lenguajes de programación brindando un primer acercamiento a la construcción básica de uno, partiendo de las definiciones de conceptos que componen un lenguaje, específicamente: sintaxis, semántica, y pragmática. Así como a la clasificación que agrupa y distingue a los diferentes lenguajes computacionales en la actualidad\footnote{Para las personas que poseen una madurez matemática superior a aquella que la que un estudiante de cuarto semestre de la Facultad de Ciencias pudiera tener se recomienda encarecidamente iniciar el estudio de este manual en el \hyperref[sec:sintax]{capítulo 3: Sintaxis}.\\ Los temas recomendados para omitir los primeros dos capítulos del manual son: estructuras discretas, juicios lógicos, estructuras recursivas, reglas de inferencia e inducción estructural. }.

\subsection*{Planteamiento}
    Este capítulo inicia con el estudio de los lenguajes de programación definiendo diferentes gramáticas para generar lenguajes pequeños y familiarizarnos con los elementos que los conforman.\\\\
    Es de particular interés analizar nuestro primer caso de estudio, el lenguaje \textsf{EAL}: Expresiones aritmético-lógicas proporcionando la gramática y esbozando como podemos interpretar las expresiones que pertenecen al mismo mediante la función de evaluación \texttt{eval}.\\\\
    Para cerrar el capítulo se discute brevemente la clasificación de los lenguajes de programación mas importantes de la actualidad basados en las características y paradigma de programación al que pertenecen. 

%    Sección 1: Composición de los lenguajes
\section{Reglas de los lenguajes}

    La lingüística tiene como propósito definir una serie de reglas que rijan la estructura fundamental que compone un lenguaje.
Las reglas que estudian dicha composición pueden ser clasificadas en sintácticas (estructura gramatical), semánticas (significado) y pragmáticas (contexto de uso).\\

    % Ejercicio 1.2
    \begin{exercise}
        Explica que es la sintaxis  de un lenguaje.     \\
        En el contexto de las ciencias de la computación, la sintaxis de un lenguaje son las reglas que definen las combinaciones de símbolos que se consideran declaraciones o expresiones correctamente estructuradas\footnote{Definición extraída de  \hyperlink{29}{[29]}, \hyperlink{30}{[30]} y \hyperlink{31}{[31]}}.
        
    \end{exercise}

\bigskip

    % Ejercicio 1.1
    \begin{exercise}
       Explica que es la semántica de un lenguaje.    \\
        En el contexto de las ciencias de la computación, la semántica describe los procesos que sigue una computadora al ejecutar un programa en un lenguaje específico. Esto se puede mostrar describiendo la relación entre la entrada y la salida de un programa, o una explicación de cómo se ejecutará el programa en una determinada plataforma, creando así un modelo de cálculo\footnote{Definición extraída de \hyperlink{20}{[20]},  \hyperlink{21}{[21]} y \hyperlink{28}{[28]}}.
    \end{exercise} 


    % Ejercicio 1.3
    \begin{exercise}
        Explica que es la pragmática de un lenguaje.   \\
         En el contexto de las ciencias de la computación la pragmática se refiere a todos aquellos elementos del lenguaje que nos permiten obtener el significado de una expresión bajo un determinado contexto.
        Un ejemplo de esto son las bibliotecas que nos permiten importar métodos previamente definidos, así la interpretación de una instrucción de esta biblioteca solo tiene un significado válido bajo el contexto que la importa\footnote{Extraído de \hyperlink{36}{[36]} y   \hyperlink{37}{[37]}}.
    \end{exercise} 

    \bigskip

%    Sección 2: Gramática de los Lenguajes.
\section{Gramáticas y semántica de expresiones}

    Los lenguajes (no solo los de programación) requieren de un conjunto de reglas que nos permita formar las expresiones que pertenecen a ellos y un método de interpretación que nos diga el significado de las expresiones formadas a partir de dicho conjunto de reglas.\\\\
    En cursos como autómatas y lenguajes formales\footnote{Conforme al plan de estudios que se imparte desde el 2013 en la Facultad de Ciencias de la Universidad Nacional Autónoma de México con clave de asignatura 1425. } se revisan este tipo de estructuras y las reglas que las generan conocidas como sintaxis del  lenguaje. Estas reglas pueden ser definidas mediante juicios lógicos o gramáticas. \\\\
    Adicionalmente se necesita una función de evaluación que nos permita obtener el valor asociado a la expresión. A esta rama se le conoce como semántica del lenguaje\footnote{
    Más adelante en este curso revisaremos a profundidad las reglas semánticas y como evaluar expresiones correctas.}.

    % Ejercicio 2.1
    \begin{exercise}
        Define la gramática para las cadenas de la forma $(ab)^*$.     \\\\
            $S$ ::= ab $S$ $|$ $\varepsilon$
    \end{exercise}

    % Ejercicio 2.2
    \begin{exercise}
        Define la gramática para números enteros.  \\\\
            $S$ ::= $D$ $|$ $D$ $S$     \\
            $D$ ::= 0 $|$1 $|$ 2 $|$ 3 $|$ 4 $|$ 5 $|$ 6 $|$ 7 $|$ 8 $|$ 9
    \end{exercise}


    %Ejercicio 2.3
    \begin{exercise}
        Define la gramática para las listas de dígitos.     \\\\
            $S$ ::= [ $DL$ ] $|$ [ $\varepsilon$ ]    \\
            $DL$ ::= $D$ $|$ $D$ , $DL$  \\
            $D$ ::= 0 $|$ 1 $|$ 2 $|$ 3 $|$ 4 $|$ 5 $|$ 6 $|$ 7 $|$ 8 $|$ 9
    \end{exercise}

    % Definición del lenguaje EAB
    \begin{definition}[]Sintaxis de {\sf Expresiones aritmético-lógicas: EAL}
        \[
            \begin{array}{rll}
                e & ::= &  n\quad |\quad {\tt true}\quad |\quad {\tt false}\quad |\quad e + e\quad |\quad e*e\quad |\quad e<e\quad 
            \end{array}
       \]
       \[
	 \begin{array}{rll}
		|\quad {\tt iszero}\;e\quad |\quad {\tt not}\;e
	\end{array}
        \]
        En donde $n\in\mathbb{N}$ y las constantes \textsf{true} y \textsf{false} son cadenas del lenguaje (nótese que no son elementos de $Bool$ aún)\footnote{Definición formulada de  \hyperlink{1}{[1]},  \hyperlink{5}{[5]}.  \hyperlink{12}{[12]},  \hyperlink{26}{[26]} y  \hyperlink{39}{[39]}.}.
    \end{definition}

    \bigskip
        
    % Definición de la función semanticá para EAB
    \begin{definition}[Semántica de \textsf{EAL}] Función de evaluación para el lenguaje de expresiones aritmético-lógicas.
        $$\llbracket\cdot\rrbracket:{\sf EAL}\to\mathbb{N} \cup \textit{Bool}$$
         Esto quiere decir que la función de evaluación recibe una expresión del lenguaje \textsf{EAL} y regresa un número natural ó una constante booleana y se define de la siguiente forma:
        \[
            \begin{array}{rll}
                \llbracket n\rrbracket&=&n\\
                \llbracket {\tt true}\rrbracket&=&{\tt True}\\
                \llbracket {\tt false}\rrbracket&=&{\tt False}\\
                \llbracket e_1 + e_2\rrbracket&=&\llbracket e_1\rrbracket+_\mathbb{N}\llbracket e_2\rrbracket\\
                \llbracket e_1*e_2\rrbracket&=&\llbracket e_1\rrbracket\times_\mathbb{N}\llbracket e_2\rrbracket\\
                \llbracket e_1<e_2\rrbracket&=&\llbracket e_1\rrbracket<_\mathbb{N}\llbracket e_2\rrbracket\\
                \llbracket {\tt iszero}\;e\rrbracket&=&{\sf zero}_\mathbb{N}\llbracket e\rrbracket\\
                \llbracket {\tt not}\;e\rrbracket&=&\lnot\llbracket e\rrbracket
            \end{array}
        \]
        En donde $ n \in \mathbb{N}$ es un natural y las constantes $\tt False$, $\tt True \in $ $Bool$.\\\\
        Los operadores $+_\mathbb{N},\times_\mathbb{N},<_\mathbb{N}$ y ${\sf zero}_\mathbb{N}$ representan las funciones de suma, producto, menor y el \textit{test} de cero definidas para los naturales respectivamente, mientras que $\lnot$ es la negación lógica para $Bool$\footnote{Definición extraída de \hyperlink{1}{[1]},  \hyperlink{5}{[5]}.  \hyperlink{12}{[12]},  \hyperlink{26}{[26]} y  \hyperlink{39}{[39]}.}.
    \end{definition}

    A partir de la definición anterior se puede implementar un intérprete para expresiones del lenguaje \textsf{EAL} mediante una función recursiva $\tt eval$ como sigue:
    \vskip 1em
    \begin{lstlisting}
        eval n          = n
        eval true       = True
        eval false      = False
        eval (e1 + e2)  = eval e1 + eval e2
        eval (e1 * e2)  = eval e1 * eval e2
        eval (e1 < e2)  = (eval e1) < (eval e2)
        eval (not e)    = not (eval e)
        eval (iszero e) = (eval e) == 0
     \end{lstlisting}

    Con esta definición ahora podemos evaluar algunas expresiones pertenecientes a nuestro lenguaje \textsf{EAL}.

    % Ejercicio 2.4
    \begin{exercise}
        Evalúa la expresión: \texttt{not} 3 $\textless$ 5 + 7  haciendo uso de la implementación del intérprete anteriormente definido para \textsf{EAL}.
        \vskip 1em
        \begin{lstlisting}
                eval(not 3 < 5 + 7)  =   	                                                                
                not(eval(3 < 5 + 7))  =          	                                                               
                not(eval(3) < eval(5 + 7))  =
                not(eval(3) <  eval(5) + eval(7)) =   
                not(3 < 5 + 7)  =                                                                            
                not(3 < 12)  =               	                                                                                                 
                not(True)  =
                False
        \end{lstlisting}
        
    \end{exercise}


    % Ejercicio 2.5
    \begin{exercise}
        Evalúa la expresión: $3$ * $\texttt{true}$ $ \textgreater$  $8$  haciendo uso de la implementación del intérprete anteriormente definido para \textsf{EAL}.
        \vskip 1em
        \begin{lstlisting}
            eval(3 * true > 8)        =                              
            eval(3 * true) > eval(8)  =            
            eval(3) * eval(true) >  8 =                                                  
            3 * True > 8              =  
            error
        \end{lstlisting}
        Expresión de \textsf{EAL} mal formada.

    \end{exercise}

    \begin{definition}
        Sea \textit{G} una gramática libre de contexto, decimos que que \textit{G} es una gramática ambigua si existe una cadena para la cual se pueda tener más de una derivación a la izquierda\footnote{Definición extraída de \hyperlink{5}{[5]}. \hyperlink{12}{[12]}, \hyperlink{40}{[40]} y \hyperlink{41}{[41]}.}.
    \end{definition}


    \begin{exercise}
        ¿La gramática definida para \textsf{EAL} es ambigua? Argumenta tu respuesta.                          \\
             Sí porque podemos derivar de dos formas distintas la cadena: $3 + 4 * 5$                            
             \begin{center}
                 \textit{e} $\rightarrow$ \\
                 \textit{e} + \textit{e} $\rightarrow$ \\
                 3 + \textit{e} $\rightarrow$ \\
                 3 + \textit{e} * \textit{e} $\rightarrow$ \\
                 3 + 4 * \textit{e} $\rightarrow$ \\
                 3 + 4 * 5.  \\
                 
                 \bigskip
                 
                 \textit{e} $\rightarrow$ \\
                 \textit{e} * \textit{e} $\rightarrow$ \\
                 \textit{e} * 5 $\rightarrow$ \\
                 \textit{e} + \textit{e} * 5 $\rightarrow$ \\
                 \textit{e} + 4 * 5 $\rightarrow$ \\
                 3 + 4 * 5.
             \end{center}
    \end{exercise}

    \bigskip

    % Ejercicio 2.7
    \begin{exercise}
        ¿Cómo eliminarías la ambigüedad de esta gramática?    \\\\                    
            Con la introducción de paréntesis para marcar la precedencia de las operaciones. 
            En el ejemplo anterior la ambigüedad queda eliminada en su totalidad:
            \begin{center}
                $(3 + (4 * 5))$  ó $((3 + 4 )* 5)$
            \end{center}
    \end{exercise}


% Sección 3: Clasificación de los lenguajes.
\section{Clasificación de los lenguajes de programación}

    En computación existen diferentes clasificaciones que nos permiten agrupar lenguajes según sus características principales, paradigma de programación al que pertenecen, tecnología involucrada en el proceso de traducción a lenguaje de máquina, etc. A continuación se enlistan algunas de estas categorías: 

    \begin{itemize}
        \item Lenguajes compilados,  que como el nombre lo indica precisan de un compilador para traducir el código de alto nivel a instrucciones de máquina que el procesador pueda ejecutar. Generalmente suelen ser mas lentos al momento de ejecutarse que el resto de las categorías  ya que requieren de un sistema sofisticado de inferencia de tipos, optimizaciones en la traducción de las instrucciones y en ocasiones entornos especiales para la ejecución de los programas como la máquina virtual de Java (\textit{JVM} por sus siglas en inglés)\footnote{Información obtenida de \hyperlink{43}{[43]}, \hyperlink{44}{[44]}, \hyperlink{45}{[45]}, \hyperlink{46}{[46]}}.\\

        \item Lenguajes interpretados, en donde las instrucciones son mapeadas a su traducción equivalente en el lenguaje de máquina cuando la línea que contiene dicha instrucción se vaya a ejecutar por el procesador\footnote{También es posible interpretar las expresiones de un lenguaje como el valor asociado a dicha expresión una vez que se evalúa. En este curso estudiaremos este proceso en el \hyperref[sec:semantics]{capítulo 4: Semántica}.}. Generalmente son menos robustos a la hora de encontrar errores de tipo, errores en la tipografía, caracteres faltantes, etc. Sin embargo este tipo de lenguajes de programación resultan particularmente útiles por que permiten el modelado y la realización de programas mucho más rápido que los lenguajes compilados\footnote{Información obtenida de \hyperlink{43}{[43]}, \hyperlink{44}{[44]}, \hyperlink{45}{[45]}, \hyperlink{46}{[46]}}.\\

        \item También se pueden clasificar por paradigma como en el caso de los lenguajes funcionales que se enfocan en la especificación de un problema mediante condiciones en lugar de la solución teniéndose la siguiente definición: 
         ''La programación funcional es un paradigma de programación en el que los programas se construyen aplicando y componiendo funciones. Íntimamente relacionado con el paradigma de programación declarativa en el que las definiciones de funciones son árboles de expresiones que asignan valores a otros valores, en lugar de una secuencia de declaraciones imperativas que actualizan el estado de ejecución del programa. ''\footnote{Definición extraída de \hyperlink{51}{[51]} y \hyperlink{52}{[52]}}\\

        \item  La programación imperativa contrasta con la clasificación anterior por su naturaleza secuencial, podemos encapsular sus características más importantes en la siguiente definición: 
         ''Un paradigma de programación de software que utiliza declaraciones que cambian el estado de un programa. De la misma manera que el modo imperativo en los lenguajes naturales expresa órdenes, un programa imperativo consta de órdenes que debe ejecutar la computadora. La programación imperativa se centra en describir cómo funciona un programa paso a paso\footnote{[\hyperlink{53}{[53]}} en lugar de descripciones de alto nivel de sus resultados esperados. ''\footnote{Extraído de \hyperlink{53}{[53]}, \hyperlink{54}{[54]}, \hyperlink{55}{[55]} y \hyperlink{56}{[56]}}\\\\
        El término se utiliza con frecuencia como un contrapunto de la programación declarativa, siendo el primero un conjunto de instrucciones para solucionar un determinado problema mientras que el segundo se centra en la especificación del problema mismo.\\

        \item     Otra clasificación que es muy socorrida en la computación es la llamada orientada a objetos que modela las entidades de un programa con su identidad, métodos, variables y que brindan características deseables como la herencia, el polimorfismo y el encapsulamiento de datos.\\\\
        Esta clasificación pertenece al paradigma imperativo dado que se trabaja estrechamente con el estado del programa en cada paso de la ejecución, sus direcciones de memoria y sus sentencias de control\footnote{Definición extraída de \hyperlink{47}{[47]}, \hyperlink{48}{[48]}, \hyperlink{49}{[49]}}.\\
    
        \item     Por último mencionaremos brevemente la clasificación de los lenguajes lógicos. Los programas que pertenecen a esta categoría se pueden definir de la siguiente forma:  ''Un programa, base de datos o base de conocimientos en un lenguaje de programación lógica es un conjunto de oraciones en forma lógica que expresan hechos y reglas sobre un dominio específico. ''\footnote{Definición extraída de \hyperlink{57}{[57]}} \\\\
        Generalmente este tipo de programas se apoyan definiendo predicados y objetos a los cuales estos predicados pueden ser aplicados así como instrucciones y operadores lógicos para dar solución a los problemas planteados en este modelo.

    
    \end{itemize}

    \bigskip

    % Ejercicio 3.3
    %\begin{exercise}
    %    Elabora un mapa mental con las características y clasificaciones que consideres más importante en los lenguajes de programación.
    %    \begin{figure}[htbp]
    %        \centerline{\includegraphics[scale=.45]{mapa.jpg}}
    %        \label{fig 1. Mapa para clasificar los lenguajes de programación.}
    %    \end{figure}
    %\end{exercise}

    %\begin{figure*}[ht]
    %    \centering
    %\end{figure*}
   %(Imagen de José M. Barranco:  ''Mapa Mental: Lenguajes de Programación) ''

    \bigskip
    \bigskip

 % Ejercicio 3.1
    \begin{exercise}
        Da una implementación del algoritmo \texttt{selectionSort} en programación imperativa.

	\begin{lstlisting}[language=C++]
	void swap(int *xp, int *yp){
	    int temp = *xp;
	    *xp = *yp;
	    *yp = temp;
	}
	 
	void selectionSort(int arr[], int n){
	    int i, j, min_idx;
	    for (i = 0; i < n-1; i++){
	        min_idx = i;
	        for (j = i+1; j < n; j++)
	        if (arr[j] < arr[min_idx])
	            min_idx = j;
	        swap(&arr[min_idx], &arr[i]);
	    }
	}
	 
	\end{lstlisting}
    \end{exercise}

      % Ejercicio 3.2
    \begin{exercise}
        Da una implementación del algoritmo \texttt{selectionSort} en programación declarativa.
	\begin{lstlisting}[language=Haskell] import Data.List (minimum, delete)
		ssort :: Ord t => [t] -> [t]
		ssort [] = []
		ssort xs = let { x = minimum xs } 
		           in  x : ssort (delete x xs)
	\end{lstlisting}
    \end{exercise}


% Ejercicios a discresión del lector   
\section{Ejercicios para el lector}

    % Ejercicio 4.1
    \begin{exercise}
        El uso del carácter especial ';' al terminar una instrucción de programa en lenguajes de programación como \textsf{C} y \textsf{Java} ¿a qué tipo de regla corresponde (sintáctica, semántica o pragmática)?
    \end{exercise}

\bigskip

    % Ejercicio 4.2
    \begin{exercise}
        Ejemplifica la aplicación de la pragmática aplicada en el contexto de los lenguajes de programación.
    \end{exercise}

\bigskip

     % Ejercicio 4.3
    \begin{exercise}
        En el contexto de los lenguajes de programación muchas veces la semántica de una expresión es explicada en lenguaje natural en la documentación del lenguaje. \\\\
        Explica por qué esto no es una semántica formal.
    \end{exercise}

\bigskip

    % Ejercicio 4.4
    \begin{exercise}
        Escribe una gramática para el alfabeto $\Sigma = \{a,b\}$ y el lenguaje
        \[L=\{a^n b^m\ \textbar\ n \textgreater 0, m \geq n \}\]
    \end{exercise}

\bigskip

     % Ejercicio 4.5
    \begin{exercise}
        Escribe una gramática para el alfabeto $\Sigma = \{a,b\}$ y el lenguaje  
        \[L= \text{cadenas palíndromas} \]
    \end{exercise}

\bigskip
    % Ejercicio 4.6
    \begin{exercise}
        Escribe una gramática para el alfabeto  $\Sigma = \{(,)\}$ y el lenguaje \[ L= \text{cadenas de paréntesis balanceados} \]
    \end{exercise}

\bigskip

    % Ejercicio 4.7
    \begin{exercise}
        Evalúa la expresión del lenguaje \textsf{EAL} utilizando la función \texttt{eval}. \\ \\ 
        $$7 * 8 < 3 * 9$$
    \end{exercise}

\bigskip

    % Ejercicio 4.8
    \begin{exercise}
        Evalúa la expresión del lenguaje \textsf{EAL} utilizando la función \texttt{eval}. \\ \\
            $$\texttt{not(iszero}( (7 * 4) + (1 * 0)))$$
    \end{exercise}

\bigskip

    % Ejercicio 4.9
    \begin{exercise}
        Escribe el algoritmo de la búsqueda binaria en un arreglo ordenado en el paradigma de programación imperativo.
    \end{exercise}
\bigskip
    % Ejercicio 4.10
    \begin{exercise}
        Escribe el algoritmo de la búsqueda binaria en un arreglo ordenado en el paradigma de programación declarativo.
    \end{exercise}

% https://opendsa-server.cs.vt.edu/OpenDSA/Books/PIFLAS21/html/IntroGrammarEx.html
% https://programming-idioms.org/idiom/124/binary-search-for-a-value-in-sorted-array/2120/haskell
% https://www.geeksforgeeks.org/binary-search/

% https://www.seas.upenn.edu/~cis1940/spring13/hw/01-intro.pdf
% https://wiki.haskell.org/Typeclassopedia
% https://github.com/JD95/haskell-problem-sets/blob/master/Lists/Problems.hs
% https://www.reddit.com/r/haskell/comments/evlmsr/beginner_exercises/

    %https://tex.stackexchange.com/questions/669517/notes-at-the-end-of-the-book-document-what-is-the-modern-way-of-doing-it